Here is a corrected proof with the two reported gaps filled explicitly.

\[
\text{For integers } m,n,\text{ write } m\mid n \text{ to mean that } n=mk
\text{ for some } k\in\mathbb Z.
\]

We first record two elementary facts that will be used in the proof.

---

### Lemma 1: A divisibility subtraction property

If \(c,x,y\in\mathbb Z\), \(c\mid x\), and \(c\mid y\), then

\[
c\mid (x-y).
\]

**Proof.** Since \(c\mid x\) and \(c\mid y\), there exist integers \(u,v\) such that

\[
x=cu,\qquad y=cv.
\]

Therefore

\[
x-y=cu-cv=c(u-v).
\]

Since \(u-v\in\mathbb Z\), we have \(c\mid (x-y)\). \(\square\)

In particular, if \(c\mid x\) and \(c\mid (x-y)\), then

\[
y=x-(x-y),
\]

so Lemma 1 implies \(c\mid y\).

---

### Lemma 2: A finite maximum property

If \(N\in\mathbb Z_{>0}\) and \(S\) is a nonempty subset of

\[
\{1,2,\dots,N\},
\]

then \(S\) has a largest element.

**Proof.** We prove this by induction on \(N\).

For \(N=1\), the only nonempty subset of \(\{1\}\) is \(\{1\}\), whose largest element is \(1\).

Assume the result is true for some \(N\ge 1\). Let \(S\) be a nonempty subset of

\[
\{1,2,\dots,N,N+1\}.
\]

If \(N+1\in S\), then \(N+1\) is the largest element of \(S\). If \(N+1\notin S\), then

\[
S\subseteq \{1,2,\dots,N\}.
\]

Since \(S\) is nonempty, the induction hypothesis implies that \(S\) has a largest element.

Thus the claim holds for \(N+1\). By induction, the lemma is proved. \(\square\)

---

### Lemma 3: If \(7\mid n^2\), then \(7\mid n\)

Let \(n\in\mathbb Z\). Suppose

\[
7\mid n^2.
\]

By the division algorithm, since \(7>0\), there exist integers \(a\) and \(r\) such that

\[
n=7a+r
\]

with

\[
r\in\{0,1,2,3,4,5,6\}.
\]

Then

\[
n^2=(7a+r)^2=49a^2+14ar+r^2.
\]

Hence

\[
n^2-r^2=49a^2+14ar=7(7a^2+2ar).
\]

Since \(7a^2+2ar\in\mathbb Z\), we have

\[
7\mid (n^2-r^2).
\]

We also assumed that \(7\mid n^2\). Therefore, by Lemma 1 applied to \(n^2\) and \(n^2-r^2\),

\[
7\mid \bigl(n^2-(n^2-r^2)\bigr).
\]

But

\[
n^2-(n^2-r^2)=r^2.
\]

Thus

\[
7\mid r^2.
\]

Now check the possible values of \(r\):

\[
\begin{array}{c|ccccccc}
r & 0 & 1 & 2 & 3 & 4 & 5 & 6 \\
\hline
r^2 & 0 & 1 & 4 & 9 & 16 & 25 & 36
\end{array}
\]

Among these numbers, the only one divisible by \(7\) is \(0\). Therefore \(r=0\). Hence

\[
n=7a,
\]

so

\[
7\mid n.
\]

This proves the lemma. \(\square\)

---

### Lemma 4: Every rational expression can be put in lowest terms

If \(p,q\in\mathbb Z\) and \(q\neq 0\), then there exist integers \(a,b\) such that

\[
\frac{p}{q}=\frac{a}{b}, \qquad b\neq 0,
\]

and \(a\) and \(b\) have no positive common divisor greater than \(1\).

**Proof.** Define

\[
D=\{d\in\mathbb Z_{>0}: d\mid p \text{ and } d\mid q\}.
\]

The set \(D\) is nonempty because \(1\mid p\) and \(1\mid q\), so

\[
1\in D.
\]

Now let \(d\in D\). Since \(d\mid q\), there exists an integer \(t\) such that

\[
q=dt.
\]

Because \(q\neq 0\), we must have \(t\neq 0\). Therefore

\[
|t|\ge 1.
\]

Since \(d>0\), we get

\[
|q|=|dt|=d|t|\ge d.
\]

Thus every \(d\in D\) satisfies

\[
1\le d\le |q|.
\]

So

\[
D\subseteq \{1,2,\dots,|q|\}.
\]

Because \(q\neq 0\), \(|q|\in\mathbb Z_{>0}\). Therefore, by Lemma 2, the nonempty set \(D\) has a largest element. Call it \(d\).

Since \(d\in D\), we have \(d\mid p\) and \(d\mid q\). Hence there exist integers \(a,b\) such that

\[
p=da,\qquad q=db.
\]

Because \(q\neq 0\), we have \(b\neq 0\). Also,

\[
\frac{p}{q}=\frac{da}{db}=\frac{a}{b}.
\]

We claim that \(a\) and \(b\) have no positive common divisor greater than \(1\). Suppose, for contradiction, that there exists an integer \(c>1\) such that

\[
c\mid a
\qquad\text{and}\qquad
c\mid b.
\]

Then there exist integers \(a_1,b_1\) such that

\[
a=ca_1,\qquad b=cb_1.
\]

Therefore

\[
p=da=dca_1,
\qquad
q=db=dcb_1.
\]

So

\[
dc\mid p
\qquad\text{and}\qquad
dc\mid q.
\]

Since \(d>0\) and \(c>1\), we have \(dc\in\mathbb Z_{>0}\) and

\[
dc>d.
\]

Thus \(dc\in D\), contradicting the fact that \(d\) was the largest element of \(D\).

Therefore \(a\) and \(b\) have no positive common divisor greater than \(1\). \(\square\)

---

### Theorem: \(\sqrt 7\) is irrational

Suppose, for contradiction, that \(\sqrt 7\) is not irrational. By the given definition, this means that there exist integers \(p,q\) with \(q\neq 0\) such that

\[
\sqrt 7=\frac{p}{q}.
\]

By Lemma 4, we may write the same number in lowest terms:

\[
\sqrt 7=\frac{a}{b},
\]

where \(a,b\in\mathbb Z\), \(b\neq 0\), and \(a\) and \(b\) have no positive common divisor greater than \(1\).

Squaring both sides gives

\[
7=\frac{a^2}{b^2}.
\]

Since \(b\neq 0\), we may multiply both sides by \(b^2\), obtaining

\[
7b^2=a^2.
\]

Thus

\[
7\mid a^2.
\]

By Lemma 3,

\[
7\mid a.
\]

So there exists an integer \(k\) such that

\[
a=7k.
\]

Substituting this into \(7b^2=a^2\), we get

\[
7b^2=(7k)^2=49k^2.
\]

Dividing both sides by \(7\), we obtain

\[
b^2=7k^2.
\]

Hence

\[
7\mid b^2.
\]

By Lemma 3 again,

\[
7\mid b.
\]

Therefore \(7\mid a\) and \(7\mid b\). So \(7\) is a positive common divisor of \(a\) and \(b\), and since

\[
7>1,
\]

this contradicts the choice of \(a\) and \(b\) as having no positive common divisor greater than \(1\).

Thus no integers \(p,q\) with \(q\neq 0\) satisfy

\[
\sqrt 7=\frac{p}{q}.
\]

Therefore, by the given definition, \(\sqrt 7\) is irrational. \(\square\)